% Hanzo Post-Quantum Cryptography: Migration Path for AI Infrastructure
\documentclass[11pt]{article}
\usepackage{shared/hanzocover}
\usepackage{amsmath, amssymb, amsthm}
\usepackage{mathtools}
\usepackage{bm}
\usepackage{graphicx}
\usepackage{booktabs}
\usepackage{hyperref}
\usepackage{enumitem}
\usepackage{xcolor}
\usepackage{float}
\usepackage{multirow}

\hypersetup{colorlinks=true,linkcolor=black,citecolor=blue,urlcolor=blue}

\newtheorem{theorem}{Theorem}
\newtheorem{lemma}[theorem]{Lemma}
\newtheorem{proposition}[theorem]{Proposition}
\newtheorem{corollary}[theorem]{Corollary}
\newtheorem{definition}{Definition}
\newtheorem{remark}{Remark}

\title{HANZO POST-QUANTUM CRYPTOGRAPHY\\[0.3em]
{\large Migration Path for\\AI Infrastructure}}
\author{
    Zach Kelling \\
    \textit{Hanzo Industries Inc (Techstars '17)} \\
    \texttt{research@hanzo.ai}
}
\date{March 2026}

\begin{document}
\hanzocoverpage

% ============================================================================
% ABSTRACT
% ============================================================================
\begin{abstract}
We describe the post-quantum cryptographic migration of Hanzo infrastructure,
replacing classical algorithms vulnerable to quantum attack with NIST-standardized
post-quantum alternatives. The migration covers three domains: (i)~ML-KEM
(Module-Lattice-Based Key Encapsulation Mechanism, formerly Kyber) for key
exchange in TLS connections between Hanzo services, client API connections, and
KMS envelope encryption, (ii)~ML-DSA (Module-Lattice-Based Digital Signature
Algorithm, formerly Dilithium) for API request signing, IAM JWT tokens, and
agent credential chains, and (iii)~SLH-DSA (Stateless Hash-Based Digital
Signature Algorithm, formerly SPHINCS+) as a conservative fallback for
long-lived root certificates. We present the hybrid transition strategy where
classical and post-quantum algorithms run in parallel during migration, the
performance impact on Hanzo Gateway request latency (1.2ms increase for ML-KEM
handshake, 0.8ms for ML-DSA verification), and the key size implications for
KMS storage and IAM token payloads. Production deployment began in January 2026
with ML-KEM for inter-service TLS; full migration completes by Q4 2026.
\end{abstract}

% ============================================================================
\section{Introduction}
\label{sec:intro}
% ============================================================================

Cryptographically relevant quantum computers (CRQCs) threaten the security
foundations of all modern infrastructure. Shor's algorithm~\cite{shor1994}
breaks RSA, ECDSA, and ECDH in polynomial time. Grover's
algorithm~\cite{grover1996} weakens symmetric ciphers and hash functions by a
quadratic factor.

Hanzo infrastructure relies on cryptography at every layer:
\begin{itemize}[leftmargin=1.1em]
    \item \textbf{TLS}: ECDHE key exchange (X25519) and ECDSA server
    certificates for all service-to-service and client-to-service connections.
    \item \textbf{IAM}: Ed25519 JWT signing for authentication tokens.
    \item \textbf{KMS}: AES-256-GCM envelope encryption with ECDH-derived
    key wrapping for secret storage.
    \item \textbf{Agent credentials}: Ed25519 certificate chains for delegated
    agent permissions.
    \item \textbf{ACI}: ECDSA (secp256k1) for blockchain transaction signing
    and compute attestations.
\end{itemize}

While large-scale CRQCs do not yet exist, the ``harvest now, decrypt later''
threat model requires action today: adversaries recording encrypted traffic can
decrypt it once quantum computers become available. Data with long
confidentiality requirements (medical records, financial data, government
communications) processed through Hanzo AI must be protected against future
quantum decryption.

\subsection{NIST Standardization}

NIST finalized three post-quantum standards in August 2024~\cite{nistpq}:
\begin{enumerate}[leftmargin=1.4em]
    \item \textbf{ML-KEM} (FIPS 203): Lattice-based key encapsulation,
    replacing ECDH for key exchange.
    \item \textbf{ML-DSA} (FIPS 204): Lattice-based digital signatures,
    replacing ECDSA/Ed25519 for signing.
    \item \textbf{SLH-DSA} (FIPS 205): Hash-based digital signatures,
    providing a conservative alternative based on minimal assumptions.
\end{enumerate}

% ============================================================================
\section{Threat Model}
\label{sec:threat}
% ============================================================================

\begin{definition}[Quantum Threat Timeline]
Let $T_Q$ be the year a CRQC capable of breaking 256-bit elliptic curve
cryptography becomes available. Let $T_D$ be the required data confidentiality
horizon. Data encrypted today is vulnerable if $T_D > T_Q$.
\end{definition}

Conservative estimates place $T_Q$ between 2030 and
2040~\cite{quantumtimeline}. For Hanzo infrastructure:
\begin{itemize}[leftmargin=1.1em]
    \item \textbf{Healthcare AI data}: $T_D = 50+$ years (HIPAA requirements).
    Migration urgency: \emph{critical}.
    \item \textbf{Financial AI data}: $T_D = 7$--$25$ years (regulatory
    retention). Migration urgency: \emph{high}.
    \item \textbf{Session tokens}: $T_D < 24$ hours (short-lived). Migration
    urgency: \emph{moderate} (but required for compliance).
    \item \textbf{Blockchain signatures}: $T_D = \infty$ (on-chain forever).
    Migration urgency: \emph{critical}.
\end{itemize}

\begin{theorem}[Migration Deadline]
For data with confidentiality horizon $T_D$ and migration timeline $T_M$
(years to complete migration), the latest safe migration start date is
$T_Q - T_M - T_D$. With $T_Q = 2033$ (median estimate), $T_M = 2$ years,
and $T_D = 10$ years: migration must begin by 2021. We are already in the
vulnerability window.
\end{theorem}

% ============================================================================
\section{ML-KEM for Key Exchange}
\label{sec:mlkem}
% ============================================================================

ML-KEM (FIPS 203) replaces ECDH key exchange. It is based on the Module
Learning with Errors (MLWE) problem.

\begin{definition}[ML-KEM]
ML-KEM consists of three algorithms:
\begin{itemize}[leftmargin=1.1em]
    \item $\texttt{KeyGen}() \to (ek, dk)$: Generate encapsulation and
    decapsulation keys.
    \item $\texttt{Encaps}(ek) \to (K, c)$: Produce shared secret $K$ and
    ciphertext $c$.
    \item $\texttt{Decaps}(dk, c) \to K$: Recover shared secret from
    ciphertext.
\end{itemize}
\end{definition}

We deploy ML-KEM-768 (NIST Security Level 3, equivalent to AES-192):

\begin{table}[H]
\centering
\small
\begin{tabular}{lrrr}
\toprule
\textbf{Parameter} & \textbf{ML-KEM-768} & \textbf{X25519} & \textbf{Ratio} \\
\midrule
Public key (bytes) & 1,184 & 32 & 37$\times$ \\
Ciphertext (bytes) & 1,088 & 32 & 34$\times$ \\
Shared secret (bytes) & 32 & 32 & 1$\times$ \\
KeyGen ($\mu$s) & 28 & 48 & 0.6$\times$ \\
Encaps ($\mu$s) & 35 & 48 & 0.7$\times$ \\
Decaps ($\mu$s) & 32 & 48 & 0.7$\times$ \\
\bottomrule
\end{tabular}
\caption{ML-KEM-768 vs X25519 parameter comparison. ML-KEM is computationally
faster but has larger key/ciphertext sizes.}
\label{tab:mlkem}
\end{table}

\subsection{Hybrid Key Exchange}

During the transition period, we use hybrid key exchange combining X25519 and
ML-KEM-768:
\begin{equation}
K_{\text{hybrid}} = \texttt{HKDF-SHA256}(K_{\text{X25519}} \| K_{\text{ML-KEM}}),
\end{equation}
ensuring security if either scheme is broken. This follows the IETF hybrid
approach~\cite{hybridtls}.

\subsection{TLS Integration}

Hanzo Gateway~\cite{hanzogateway} implements hybrid TLS 1.3 with the
\texttt{X25519MLKEM768} named group:
\begin{enumerate}[leftmargin=1.4em]
    \item \texttt{ClientHello} includes both X25519 and ML-KEM-768 key shares
    (total 1,248 bytes vs.\ 32 bytes for X25519 alone).
    \item \texttt{ServerHello} includes both encapsulated keys.
    \item The handshake shared secret combines both derivations.
    \item Application data is encrypted with AES-256-GCM (quantum-safe at
    128-bit security via Grover's bound).
\end{enumerate}

\begin{remark}
The larger ClientHello (1,248 vs.\ 32 bytes) fits within a single TCP segment
and does not require an additional round trip. Measured TLS handshake overhead:
1.2ms additional latency (from 3.1ms to 4.3ms for a cold connection).
\end{remark}

\subsection{KMS Key Wrapping}

Hanzo KMS~\cite{hanzokms} envelope encryption currently wraps data encryption
keys (DEKs) using ECDH-derived keys. The post-quantum migration replaces the
key wrapping:
\begin{enumerate}[leftmargin=1.4em]
    \item Master KEK is derived via ML-KEM encapsulation.
    \item DEK encryption uses AES-256-GCM with the ML-KEM-derived key.
    \item Existing secrets are re-wrapped during the next rotation cycle
    (within 90 days per KMS rotation policy).
\end{enumerate}

% ============================================================================
\section{ML-DSA for Digital Signatures}
\label{sec:mldsa}
% ============================================================================

ML-DSA (FIPS 204) replaces Ed25519 and ECDSA for digital signatures. It is
based on the Module Learning with Errors and Module Short Integer Solution
(MLWE/MSIS) problems.

\begin{definition}[ML-DSA]
ML-DSA consists of three algorithms:
\begin{itemize}[leftmargin=1.1em]
    \item $\texttt{KeyGen}() \to (vk, sk)$: Generate verification and
    signing keys.
    \item $\texttt{Sign}(sk, m) \to \sigma$: Produce signature $\sigma$ on
    message $m$.
    \item $\texttt{Verify}(vk, m, \sigma) \to \{0, 1\}$: Verify signature
    validity.
\end{itemize}
\end{definition}

We deploy ML-DSA-65 (NIST Security Level 3):

\begin{table}[H]
\centering
\small
\begin{tabular}{lrrr}
\toprule
\textbf{Parameter} & \textbf{ML-DSA-65} & \textbf{Ed25519} & \textbf{Ratio} \\
\midrule
Public key (bytes) & 1,952 & 32 & 61$\times$ \\
Signature (bytes) & 3,309 & 64 & 52$\times$ \\
Sign ($\mu$s) & 120 & 22 & 5.5$\times$ \\
Verify ($\mu$s) & 45 & 18 & 2.5$\times$ \\
\bottomrule
\end{tabular}
\caption{ML-DSA-65 vs Ed25519 parameter comparison.}
\label{tab:mldsa}
\end{table}

\subsection{IAM JWT Migration}

Hanzo IAM~\cite{hanzoiam} JWT tokens currently use Ed25519 signatures.
The migration proceeds in three phases:

\begin{enumerate}[leftmargin=1.4em]
    \item \textbf{Phase 1 (current)}: IAM signs with both Ed25519 and ML-DSA-65.
    Tokens carry dual signatures. Verifiers accept either.
    \item \textbf{Phase 2 (Q3 2026)}: New tokens are signed with ML-DSA-65 only.
    Verifiers accept both algorithms (backward compatibility for tokens issued
    in Phase 1).
    \item \textbf{Phase 3 (Q4 2026)}: Ed25519 verification is removed. All
    tokens are ML-DSA-65.
\end{enumerate}

\begin{remark}
The larger ML-DSA-65 signature (3,309 bytes vs.\ 64 bytes) increases JWT size
from $\sim$500 bytes to $\sim$3,800 bytes. This remains well within HTTP header
limits and does not impact API performance measurably.
\end{remark}

\subsection{API Request Signing}

External API requests to Hanzo Gateway use HMAC-SHA256 for authentication
(quantum-safe at 128-bit security). Internal service-to-service requests use
mTLS with ML-DSA-65 certificates.

\subsection{Agent Credential Chains}

Agent credentials~\cite{hanzoagentsdk} use certificate chains signed with
Ed25519. Migration replaces the signature algorithm while preserving the chain
structure:
\begin{equation}
\text{Root}_{ML\text{-}DSA} \to \text{Org}_{ML\text{-}DSA} \to
\text{Agent}_{ML\text{-}DSA}.
\end{equation}

During the hybrid period, cross-signed certificates bridge old and new root
keys.

% ============================================================================
\section{SLH-DSA for Root Certificates}
\label{sec:slhdsa}
% ============================================================================

SLH-DSA (FIPS 205) provides a conservative backup for long-lived root
certificates. Unlike ML-DSA (lattice-based), SLH-DSA security relies only on
hash function collision resistance---an assumption with decades of confidence.

\begin{definition}[SLH-DSA Security Basis]
SLH-DSA is secure if the underlying hash function (SHA-256 or SHAKE-256) is
second-preimage resistant, a property believed to hold even against quantum
adversaries (Grover reduces security by at most a quadratic factor).
\end{definition}

We use SLH-DSA-SHA2-192s (NIST Level 3, ``small'' variant) for root
certificates:

\begin{table}[H]
\centering
\small
\begin{tabular}{lrr}
\toprule
\textbf{Parameter} & \textbf{SLH-DSA-192s} & \textbf{ML-DSA-65} \\
\midrule
Public key (bytes) & 48 & 1,952 \\
Signature (bytes) & 16,224 & 3,309 \\
Sign (ms) & 180 & 0.12 \\
Verify (ms) & 4.2 & 0.045 \\
\bottomrule
\end{tabular}
\caption{SLH-DSA-192s vs ML-DSA-65. SLH-DSA has smaller keys but much larger
signatures and slower operations.}
\label{tab:slhdsa}
\end{table}

Root certificates are verified rarely (once per TLS session establishment, then
cached), so the 4.2ms verification overhead is acceptable.

% ============================================================================
\section{ACI and Blockchain Migration}
\label{sec:blockchain}
% ============================================================================

Blockchain signatures present unique challenges: transactions signed today are
verified forever. A quantum adversary can forge historical transactions if the
signing algorithm is broken.

\subsection{Lux Network Coordination}

The Lux blockchain~\cite{luxwhitepaper} is independently migrating to
post-quantum signatures. Hanzo ACI coordinates with Lux for:
\begin{itemize}[leftmargin=1.1em]
    \item \textbf{Consensus signatures}: Validator attestations migrate from
    ECDSA to ML-DSA in a coordinated network upgrade.
    \item \textbf{Transaction signatures}: New transaction types support
    ML-DSA-65 signatures alongside legacy ECDSA (secp256k1).
    \item \textbf{Address derivation}: Post-quantum addresses are derived from
    ML-DSA-65 public keys, distinct from legacy ECDSA addresses.
\end{itemize}

\subsection{Threshold Signing Migration}

Hanzo Threshold Signing~\cite{hanzothreshold} migrates from FROST (Ed25519)
and CGGMP21 (ECDSA) to post-quantum threshold schemes:
\begin{itemize}[leftmargin=1.1em]
    \item \textbf{Lattice-based threshold signatures}: Research prototypes of
    threshold ML-DSA exist but are not yet production-ready. Expected
    standardization: 2027.
    \item \textbf{Interim approach}: Hash-based threshold signatures using
    SLH-DSA with Shamir secret sharing of the hash-based signing key. This
    is less efficient but provably secure.
\end{itemize}

% ============================================================================
\section{Performance Impact}
\label{sec:performance}
% ============================================================================

\begin{table}[H]
\centering
\small
\begin{tabular}{lrrr}
\toprule
\textbf{Operation} & \textbf{Classical} & \textbf{PQ Hybrid} & \textbf{Overhead} \\
\midrule
TLS handshake (cold) & 3.1ms & 4.3ms & +1.2ms \\
TLS handshake (resumed) & 0.8ms & 0.8ms & +0.0ms \\
JWT signing & 0.022ms & 0.142ms & +0.120ms \\
JWT verification & 0.018ms & 0.063ms & +0.045ms \\
API request (end-to-end) & 12ms & 13.2ms & +1.2ms \\
KMS secret fetch & 5ms & 5.2ms & +0.2ms \\
\bottomrule
\end{tabular}
\caption{Performance impact of post-quantum migration. Overhead is within
acceptable bounds for all operations.}
\label{tab:perf}
\end{table}

The total API request overhead of 1.2ms (from 12ms to 13.2ms) represents a
10\% increase, dominated by the TLS handshake on cold connections. Persistent
connections (used for all service-to-service communication) see negligible
overhead.

% ============================================================================
\section{Key Size and Storage Impact}
\label{sec:storage}
% ============================================================================

Post-quantum key sizes significantly exceed classical equivalents. Impact on
Hanzo systems:
\begin{itemize}[leftmargin=1.1em]
    \item \textbf{KMS}: Each ML-KEM-768 public key stored alongside the
    classical X25519 key adds 1,184 bytes per secret. For 12,000 secrets:
    14MB additional storage (negligible).
    \item \textbf{IAM JWKS endpoint}: ML-DSA-65 public keys in the JWKS
    response increase payload from 200 bytes to 2.2KB. Cached by verifiers.
    \item \textbf{JWT tokens}: Dual-signed tokens (Ed25519 + ML-DSA-65) grow
    from $\sim$500 bytes to $\sim$4.3KB. Within HTTP header limits.
    \item \textbf{Agent credential chains}: 3-level ML-DSA-65 chains total
    $\sim$16KB (3 signatures + 3 public keys). Cached per agent session.
\end{itemize}

% ============================================================================
\section{Related Work}
\label{sec:related}
% ============================================================================

\textbf{Lux post-quantum}~\cite{luxwhitepaper} implements post-quantum
cryptography at the consensus and transaction layer. Hanzo coordinates its
infrastructure migration with Lux to ensure compatibility.

\textbf{Google Chrome}~\cite{chromepq} deployed X25519MLKEM768 in TLS 1.3
starting in Chrome 124 (2024). Our TLS configuration is compatible with
Chrome's implementation.

\textbf{Signal}~\cite{signalpq} adopted the PQXDH protocol combining X25519
and ML-KEM-768 for post-quantum-secure key exchange. Our approach follows the
same hybrid pattern.

\textbf{NIST PQC}~\cite{nistpq} standardized ML-KEM, ML-DSA, and SLH-DSA in
August 2024. We adopt all three standards at their Level 3 parameter sets.

% ============================================================================
\section{Conclusion}
\label{sec:conclusion}
% ============================================================================

Hanzo post-quantum migration replaces vulnerable classical cryptography across
all infrastructure: ML-KEM-768 for key exchange (TLS, KMS), ML-DSA-65 for
digital signatures (IAM, agent credentials, API signing), and SLH-DSA-192s for
root certificates. The hybrid transition strategy ensures security if either
classical or post-quantum schemes are broken. Performance overhead is bounded
at 1.2ms for cold TLS handshakes and 0.045ms for signature verification.
Migration began in January 2026 with inter-service TLS and completes by Q4
2026 with full ML-DSA JWT signing.

% ============================================================================
% BIBLIOGRAPHY
% ============================================================================
\begin{thebibliography}{15}

\bibitem{hanzogateway}
{Hanzo AI Research}.
\newblock Hanzo Gateway: {API} gateway with {KrakenD} backend.
\newblock \emph{Hanzo AI Research}, 2024.

\bibitem{hanzoiam}
{Hanzo AI Research}.
\newblock Hanzo {IAM}: Identity and access management for {AI} platforms.
\newblock \emph{Hanzo AI Research}, 2024.

\bibitem{hanzokms}
{Hanzo AI Research}.
\newblock Hanzo {KMS}: Secrets management with per-environment rotation and audit trail.
\newblock \emph{Hanzo AI Research}, September 2021.

\bibitem{hanzoagentsdk}
{Hanzo AI Research}.
\newblock Hanzo Agent {SDK}: Framework for building {AI} agents.
\newblock \emph{Hanzo AI Research}, 2025.

\bibitem{hanzothreshold}
{Hanzo AI Research}.
\newblock Hanzo Threshold Signing: {MPC} key management for {AI} infrastructure.
\newblock \emph{Hanzo AI Research}, April 2026.

\bibitem{luxwhitepaper}
{Lux Partners}.
\newblock Lux: A scalable multi-consensus blockchain with post-quantum cryptography.
\newblock \emph{Lux Network Technical Report}, 2024.

\bibitem{nistpq}
{NIST}.
\newblock Post-quantum cryptography standardization.
\newblock \emph{NIST FIPS 203, 204, 205}, August 2024.

\bibitem{shor1994}
P.~Shor.
\newblock Algorithms for quantum computation: Discrete logarithms and factoring.
\newblock In \emph{Proc.\ FOCS}, 1994.

\bibitem{grover1996}
L.~Grover.
\newblock A fast quantum mechanical algorithm for database search.
\newblock In \emph{Proc.\ STOC}, 1996.

\bibitem{quantumtimeline}
M.~Mosca.
\newblock Cybersecurity in an era with quantum computers: Will we be ready?
\newblock \emph{IEEE Security \& Privacy}, 16(5):38--41, 2018.

\bibitem{hybridtls}
D.~Stebila, S.~Fluhrer, and S.~Gueron.
\newblock Hybrid key exchange in {TLS 1.3}.
\newblock \emph{IETF Internet-Draft}, 2024.

\bibitem{chromepq}
{Google Chrome Team}.
\newblock Post-quantum cryptography in {Chrome}.
\newblock \emph{Chromium Blog}, 2024.

\bibitem{signalpq}
{Signal Foundation}.
\newblock {PQXDH}: Post-quantum extended {Diffie-Hellman}.
\newblock \emph{Signal Technical Documentation}, 2023.

\end{thebibliography}

\end{document}
